\section{Reduced HNN adjunctions}\label{sec:hnn}

We now construct the extension that will join two pure-state classes.
The construction itself has a more general input: two trace-preserving
isomorphic subalgebras of a faithfully tracial algebra.  It adjoins a
unitary implementing the isomorphism, while prescribing the expectation
of every word in that unitary.  These prescriptions will also let us
restrict the construction to smaller bases and interchange two
adjunctions with data in the same base.

The matrix-corner construction below is the reduced HNN construction of
Ueda; see \cite[Proposition~2.2 and Remark~2.3]{UedaPreprint2008} and
\cite{Ueda2008}.  We give the details, including the faithful
representations and expectations needed later.  The corresponding
expectation and embedding results for reduced amalgamated free products
appear in \cite[Lemma~1.1 and Theorem~1.3]{BlanchardDykema2001}.

\subsection{The reduced adjunction}

Let $(A,\tau)$ be a unital $C^*$-algebra with faithful tracial state.
Fix unital subalgebras $D_1,D_2\subset A$, trace-preserving conditional
expectations $E_i:A\to D_i$, and a trace-preserving isomorphism
$\theta:D_1\to D_2$.  A word in $A$ and a unitary $v$ has the form
\begin{equation}\label{na:hnn-word}
 a_0v^{\epsilon_1}a_1\cdots v^{\epsilon_m}a_m,
 \qquad a_t\in A,\quad\epsilon_t\in\{+1,-1\}.
\end{equation}
For $m\geq1$, call it \emph{reduced} if every internal coefficient
satisfies
\begin{equation}\label{na:hnn-turns}
 \begin{cases}
 E_1(a_t)=0,&(\epsilon_t,\epsilon_{t+1})=(+1,-1),\\
 E_2(a_t)=0,&(\epsilon_t,\epsilon_{t+1})=(-1,+1).
 \end{cases}
\end{equation}
There is no condition at equal consecutive signs.

\begin{proposition}\label{na:hnn}
Let $(A,\tau)$ be a unital $C^*$-algebra with faithful tracial state,
let $D_1,D_2\subset A$ be unital subalgebras with trace-preserving
conditional expectations $E_i:A\to D_i$, and let
$\theta:D_1\to D_2$ be a trace-preserving isomorphism.  There is a
unital extension with a faithful tracial state extending $\tau$,
\[
 H(A)=C^*(A,v),\qquad vdv^*=\theta(d)\quad(d\in D_1),
\]
where $v$ is called the \emph{stable unitary}, and a trace-preserving
conditional expectation $E_A:H(A)\to A$
annihilating every nonempty reduced word
\eqref{na:hnn-word}.  Among faithfully tracial extensions generated by
$A$ and such a unitary, the relation and vanishing of the trace on
nonempty reduced words determine $H(A)$ uniquely, by an isomorphism
fixing $A$ and $v$.
\end{proposition}

We construct this adjunction as a matrix corner of an amalgamated free product.
The proof below records the module and its two compression expectations;
these will also be used in the pure-state bridge of Section~\ref{sec:bridge}.

\subsection{Construction from an amalgamated free product}

An element is \emph{centered} for a conditional expectation $E$ if it
belongs to $\ker E$.

\begin{lemma}\label{na:free-product}
Let $F$ and $G$ be unital $C^*$-algebras with faithful tracial states,
containing a common unital subalgebra $D$ on which the traces agree.
Suppose that $E_D^F:F\to D$ and $E_D^G:G\to D$ are trace-preserving
conditional expectations.  There is a unital $C^*$-algebra $Z$ generated
by faithful copies of $F$ and $G$ agreeing on $D$, with a faithful
tracial state extending both traces and faithful trace-preserving conditional
expectations
\[
 E_D:Z\longrightarrow D,
 \qquad E_F:Z\longrightarrow F,
 \qquad E_D=E_D^F\circ E_F.
\]
The map $E_D$ vanishes on every nonempty alternating product of elements
of $\ker E_D^F$ and $\ker E_D^G$.  The map $E_F$ vanishes on every product
\begin{equation}\label{na:free-core}
 a_0 b_1 a_1\cdots b_m a_m,
 \qquad m\geq1,
\end{equation}
where $a_0,a_m\in F$, $b_t\in\ker E_D^G$, and
$a_t\in\ker E_D^F$ for $0<t<m$.  The linear span of $F$ and these
products is dense in $Z$.
\end{lemma}

\begin{proof}
For $i\in\{F,G\}$, complete $i$ as a right Hilbert $D$-module $L_i$
with inner product
\[
 \langle x,y\rangle_D=E_D^i(x^*y).
\]
Left multiplication by $a\in i$ is bounded, since
\[
 E_D^i(x^*a^*ax)\leq\|a\|^2E_D^i(x^*x),
\]
and its adjoint is left multiplication by $a^*$.  The expectation gives
an orthogonal decomposition
\[
 L_i=D\oplus L_i^\circ,
 \qquad L_i^\circ=\overline{\ker E_D^i}.
\]
Both summands are invariant under left and right multiplication by $D$.
Form the Hilbert $D$-module
\begin{equation}\label{na:free-module}
 K=D\oplus
 \bigoplus_{\substack{m\geq1,\ i_t\in\{F,G\}\\i_t\ne i_{t+1}}}
 L_{i_1}^\circ\otimes_D\cdots\otimes_D L_{i_m}^\circ.
\end{equation}
Here $\otimes_D$ is the interior tensor product of Hilbert modules; the
left $D$-action on a word acts on its first factor.  The first summand
$D$ is called the vacuum summand, with distinguished vector $1$.

Let $K^{(i)}$ be the vacuum summand together with all words not starting
with $i$.  The decomposition of $L_i$ identifies
\[
 K=L_i\otimes_D K^{(i)}.
\]
Represent $i$ by left multiplication on the first tensor factor.  The
two representations agree on $D$.  They are faithful: an element $a$
annihilating the vacuum satisfies $E_D^i(a^*a)=0$, so its faithful
factor trace gives $a=0$.  Let $Z$ be the $C^*$-algebra they generate in
the adjointable operators on $K$.

A centered element of $i$ acting on $K^{(i)}$ adds its own tensor
factor at the front.  Thus a nonempty alternating centered product has
zero vacuum compression.  Centering arbitrary factor elements shows
that $D$ and these products span a dense subspace of $Z$.  Vacuum
compression consequently takes values in the left copy of $D$ and
defines the conditional expectation $E_D$.

Write $\tau_D$ for the common trace on $D$.  We next prove that
$\tau_Z=\tau_D\circ E_D$ is tracial and faithful.  For traciality,
compare $\tau_Z(aw)$ and $\tau_Z(wa)$ when $a$ belongs to one factor
and $w$ is an alternating centered word.  If $w$ has length at least
two, both values vanish.  Indeed, multiplication at a same-factor end,
followed by centering that end product, leaves either an alternating
centered word or a $D$-coefficient times a nonempty such word.  At a
different-factor end, first center $a$ and use $D$-bimodularity.  For a
one-letter word, use the factor trace when the factors agree; when they
differ, both values are zero.  For $w\in D$, use the factor trace.
Density and repeated cyclic movement of a factor now prove traciality
on all of $Z$.

Suppose $\tau_Z(z^*z)=0$.  For every $t\in Z$, traciality gives
\[
 0\leq\tau_Z(t^*z^*zt)
 =\tau_Z(ztt^*z^*)
 \leq\|t\|^2\tau_Z(zz^*)=0.
\]
The trace on $D$ is faithful, so $E_D((zt)^*(zt))=0$.  This is the
$D$-valued inner product of $zt1$ with itself, and hence $zt1=0$.  Products of
centered letters create all elementary tensors in
\eqref{na:free-module}; thus $Z1$ is dense in $K$.  It follows that
$z=0$, proving faithfulness of $\tau_Z$.

Finally, let $P_F$ be the orthogonal projection onto
$D\oplus L_F^\circ=L_F$.  This submodule reduces the $F$-action.
A core $w=b_1a_1\cdots a_{m-1}b_m$ as in
\eqref{na:free-core} sends $L_F$ into words starting with $G$, so
$P_FwP_F=0$.  Since $P_F$ commutes with $F$, arbitrary endpoint
coefficients satisfy
\[
 P_Fa_0wa_mP_F=a_0P_FwP_Fa_m=0.
\]
Successive centering shows that $F$ and the products
\eqref{na:free-core} span a dense subspace.  Their compressions belong
to the faithful left copy of $F$ on $L_F$, which is norm closed.
Compression therefore defines a conditional expectation $E_F:Z\to F$.
The vacuum lies in $L_F$, giving $E_D=E_D^F\circ E_F$ and trace
preservation.  Faithfulness follows from faithfulness of $\tau_Z$.
\end{proof}

\begin{proof}[Proof of Proposition~\ref{na:hnn}]
Take
\[
 F=M_2(A),\qquad G=M_2(D_1),\qquad D=D_1\oplus D_2.
\]
Embed $D$ in $F$ by $(d_1,d_2)\mapsto\operatorname{diag}(d_1,d_2)$
and in $G$ by
$(d_1,d_2)\mapsto\operatorname{diag}(d_1,\theta^{-1}(d_2))$.
The normalized matrix traces agree on $D$, and the maps
\[
 E_D^F(a)=(E_1(a_{11}),E_2(a_{22})),\qquad
 E_D^G(b)=(b_{11},\theta(b_{22}))
\]
are trace-preserving conditional expectations.  Apply
Lemma~\ref{na:free-product} to obtain $Z$, and retain the module
notation used in its proof.
The diagonal matrix projections $e_i=e_{ii}$ are common to the two
factors.  Set
\begin{equation}\label{na:hnn-corner}
 s=e_{21}^G,\qquad
 H(A)=e_1Ze_1,\qquad
 v=e_{12}^Fs,
\end{equation}
and identify $A$ with $e_1Fe_1$.  Then $v^*v=vv^*=e_1$.
For $d\in D_1$, the element $sds^*$ is the second-corner copy of
$\theta(d)$ in $F$, so $vdv^*=\theta(d)$ in the first corner.

The matrix units of $F$ identify $Z$ with $M_2(e_1Ze_1)$.  Under this
identification, the matrix algebra over $C^*(A,v)$ contains $F$ and
$s=e_{21}^Fv$.  It contains $G$ as well, because $G$ is generated by
$D$ and $s,s^*$.  Hence $H(A)=C^*(A,v)$.  The trace on the corner is
$2\tau_Z$, and its restriction to $A$ is $\tau$.  The restriction of
$E_F$ gives the required faithful trace-preserving expectation onto
$A$.

To check reduced words, write $q=e_{12}^F$, so $v=qs$ and
$v^*=s^*q^*$.  Between consecutive $G$-letters the internal
$F$-coefficients are
\[
\begin{array}{c|cccc}
 (\epsilon_t,\epsilon_{t+1}) &(+,+)&(+,-)&(-,+)&(-,-)\\ \hline
 \text{coefficient} &a_tq&a_t&q^*a_tq&q^*a_t.
\end{array}
\]
The first and last are off diagonal.  The middle two have zero
$D$-expectation precisely under the respective conditions in
\eqref{na:hnn-turns}.  Since $s,s^*$ are centered, the expanded word
is a product of the form \eqref{na:free-core}, and its expectation
onto $F$, hence onto $A$, is zero.

For uniqueness, split a coefficient at an opposite-sign turn as
$(a-E_i(a))+E_i(a)$.  The first term is centered at that turn, while
the second cancels two stable letters using the unitary relation.
Repeating this operation expresses every polynomial as an $A$-term
plus nonempty reduced words.  Thus the stipulated data determine the
trace of every polynomial, and in particular all moments of $x^*x$
for each polynomial $x$.

For a faithful tracial state $t$ on any $C^*$-algebra,
\begin{equation}\label{na:moment-norm}
 \|x\|=\lim_{r\to\infty}t((x^*x)^r)^{1/(2r)}.
\end{equation}
To see this, restrict $t$ to the commutative algebra generated by
$x^*x$ and the unit.  Its spectral measure has full support, by
faithfulness, and its $r$th-root moments tend to the spectral maximum.
Consequently identical polynomial trace moments give identical
polynomial norms.  The generator correspondence is therefore
isometric and extends to the asserted isomorphism.  This argument
requires no expectation in the comparison extension beyond the stated
trace-vanishing rule.
\end{proof}

We will also use the explicit spanning description of the matrix
algebra $Z$.  In the model \eqref{na:hnn-corner}, every centered
$G$-element is a sum of $D$-multiples of $s$ and $s^*$.  Absorb these
$D$-coefficients into neighboring $F$-coefficients; their internal
centering is preserved by bimodularity.  Thus $F$ and the products
\begin{equation}\label{na:matrix-word-span}
 a_0b_1a_1\cdots b_ma_m,
 \quad m\geq1,\quad b_t\in\{s,s^*\},\quad a_t\in F,
 \quad E_D^F(a_t)=0\ (0<t<m)
\end{equation}
have dense linear span in $Z$.

\begin{corollary}\label{na:hnn-restriction}
With the hypotheses of Proposition~\ref{na:hnn}, let $A'\subset A$
be a unital subalgebra containing $D_1$ and $D_2$.  Then
$C^*(A',v)\subset H(A)$ is the reduced adjunction $H(A')$ formed with
the restricted trace and expectations $E_i|_{A'}$.  The expectation
$E_A$ restricts to the expectation $H(A')\to A'$.
\end{corollary}

\begin{proof}
Because $D_i\subset A'$, the restrictions $E_i|_{A'}$ remain
conditional expectations onto $D_i$.  The subalgebra $C^*(A',v)$
inherits a faithful trace, the unitary relation and the reduced-word
trace vanishing.  Proposition~\ref{na:hnn} identifies it with $H(A')$.
The reduction of polynomials in the preceding proof shows that $E_A$
maps this subalgebra into $A'$ and has the stated restriction.
\end{proof}

\subsection{Interchanging adjunctions and taking limits}

The construction now has its two required analytic features: a
faithful expectation onto the base and a complete prescription of
reduced-word moments.  The latter makes it possible to change the
order of two adjunctions whose data were already present in a common
base.

\begin{proposition}\label{na:reorder}
Let $O$ be a unital faithfully tracial $C^*$-algebra.  For each
$t\in\{j,k\}$, let $D_1^t,D_2^t\subset O$, $E_i^t:O\to D_i^t$
and $\theta_t:D_1^t\to D_2^t$ satisfy the hypotheses of
Proposition~\ref{na:hnn}.  When forming an adjunction over $H_j(O)$
or $H_k(O)$, extend the other data expectations by composition with
the expectation onto $O$.  There is a trace-preserving isomorphism
\begin{equation}\label{na:reorder-equation}
 H_k(H_j(O))\ \cong\ H_j(H_k(O))
\end{equation}
fixing $O$ and the named stable unitaries $v_j,v_k$.
\end{proposition}

\begin{proof}
Consider $C=H_k(H_j(O))$.  Corollary~\ref{na:hnn-restriction} places
both $H_j(O)$ and $H_k(O)$ inside $C$.  Compose the expectations
$C\to H_j(O)\to O$, and call the result $E_O$.  Its restriction to
each smaller adjunction is that adjunction's expectation onto $O$,
as can be checked on its reduced words.

We prove that the two smaller adjunctions are free over $O$: an
alternating product of elements from their respective kernels of
$E_O$ has zero $E_O$-expectation.  Approximate each centered
$H_k(O)$-element by finite sums of nonempty reduced $k$-words over
$O$.  Such approximations follow by reducing a polynomial
approximation and subtracting its $O$-expectation.  Keep the centered
$H_j(O)$-elements as coefficients.

When two $k$-words are joined across a centered element
$b\in H_j(O)$, the new internal coefficient has the form $aba'$,
where $a,a'\in O$.  Its expectation onto $O$ is
$aE_O(b)a'=0$, so its expectation onto either $D_i^k$ is also zero.
The resulting word is therefore a nonempty reduced $k$-word over
$H_j(O)$, including at every new opposite-sign turn.  Initial and
final $H_j(O)$-coefficients are allowed.  Its expectation onto
$H_j(O)$ vanishes.  Passing to norm limits proves the freeness
assertion.  An alternating product with no $k$-block is a single
centered $j$-element and is already killed by $E_O$.

The individual expectations and this freeness determine all joint
$O$-valued moments.  Indeed, center the factors of a joint word; the
fully alternating centered term vanishes, and every other term has
fewer blocks after its $O$-coefficients are absorbed into neighbors.
The resulting recursion is symmetric in $j$ and $k$.  Applying the
same argument in the reverse construction gives equal scalar trace
moments for every polynomial in $O,v_j,v_k$.  Faithfulness and
\eqref{na:moment-norm} give the desired isomorphism.
\end{proof}

By Lemma~\ref{na:trace-expectation}, restrictions and compositions of
these expectations agree whenever they have the same range and
preserve the specified trace.

\begin{lemma}\label{na:hnn-limit}
Let $\lambda$ be a nonzero limit ordinal, and let
$(O_\beta)_{\beta<\lambda}$ be an increasing family of unital
$C^*$-algebras with compatible faithful tracial states.  Suppose that
fixed subalgebras $D_1,D_2$ lie in every $O_\beta$, with a fixed
trace-preserving isomorphism $\theta:D_1\to D_2$ and compatible
trace-preserving conditional expectations $E_i^\beta:O_\beta\to D_i$.
Use the embeddings of Corollary~\ref{na:hnn-restriction} to identify
the reduced adjunctions $H(O_\beta)$ with the same stable unitary $v$.
Then, with
\[
 O=\overline{\bigcup_{\beta<\lambda}O_\beta},\qquad
 R=\overline{\bigcup_{\beta<\lambda}H(O_\beta)},
\]
the expectations extend to $O\to D_i$, and $R=H(O)$ with its
specified generators and trace.
\end{lemma}

\begin{proof}
The compatible norm-one expectations extend continuously to $O$ and
retain complete positivity, bimodularity and trace preservation.
Denote these extensions by $E_i$.
The compatible traces extend faithfully to both completions.  For
example, the quotient of $R$ by its trace-null ideal is injective,
hence isometric, on every earlier faithfully tracial algebra.  It is
therefore isometric on their dense union and on $R$, so that ideal
is zero.  The same argument applies to $O$.

We have $R=C^*(O,v)$, and the unitary relation persists.  Consider a
nonempty reduced word with coefficients in $O$.  Approximate its
finitely many coefficients in one common earlier $O_\beta$.  At a
turn requiring $E_i(a)=0$, replace an approximation $b$ by
$b-E_i^\beta(b)$.  This remains in $O_\beta$, is exactly centered,
and differs from $a$ by at most twice the original error.  Thus the
given word is a norm limit of reduced words at earlier stages, all
with trace zero.  Proposition~\ref{na:hnn} identifies $R$ with $H(O)$.
This finite-coefficient argument requires no assumption on the
cofinality of $\lambda$.
\end{proof}
